% Erzeugt aus den Markdown-Bögen des Werkzeug-Repos. Nicht von Hand ändern.
\subsection*{Kopf des Vorher-Bogens}

\begin{quote}
Sag die Antwort, die Sitzungsleitung setzt das Kreuz. Sicherheit: 1 heißt geraten, 5 heißt sicher. Aus dem Kopf antworten. „Weiß ich nicht“ ist immer eine gute Antwort, wenn du es nicht weißt. Der Code in dieser Sitzung ist ein Stand vom September 2024, vor späteren Umbauten. Antworte für diesen Stand; wenn du ihn nicht kennst, sag „weiß ich nicht“.

\end{quote}

\subsection*{Kopf des Nachher-Bogens}

\begin{quote}
Sag die Antwort, die Sitzungsleitung setzt das Kreuz. Sicherheit: 1 heißt geraten, 5 heißt sicher. Aus dem Kopf antworten, nicht nachsehen. „Weiß ich nicht“ ist weiter eine gute Antwort. Der Code in dieser Sitzung ist ein Stand vom September 2024, vor späteren Umbauten. Antworte für diesen Stand.

\end{quote}

\subsection*{Teil B: Zusammenhänge}

\noindent\begin{minipage}{\linewidth}
\textbf{Frage 1}\par
Woher bekommt die Tabelle die Liste der Studierenden?

\begin{itemize}[nosep,leftmargin=2.4em]
\item[\kreuz] a\quad Die Tabelle selbst holt die Liste direkt vom Server
\item[\kreuz] b\quad Ein Baustein um die Tabelle herum holt und verteilt sie
\item[\kreuz] c\quad Die Seite holt die Liste und gibt sie der Tabelle mit
\item[\kreuz] d\quad Weiß ich nicht
\end{itemize}
\noindent Sicherheit (1 bis 5):~\rule{2.5em}{0.4pt}\par
\end{minipage}

\medskip

\noindent\begin{minipage}{\linewidth}
\textbf{Frage 2}\par
Im Ausgangsstand des Codes: Wo liegt die Aktion, die der Knopf „Import students“ auslöst?

\begin{itemize}[nosep,leftmargin=2.4em]
\item[\kreuz] a\quad In der Kopfzeile der Seite, außerhalb des Bausteins, der die Liste holt
\item[\kreuz] b\quad In der Leiste über der Tabelle, neben der Suche nach Studierenden
\item[\kreuz] c\quad Innerhalb des Bausteins, der die Liste holt und an die Tabelle gibt
\item[\kreuz] d\quad Weiß ich nicht
\end{itemize}
\noindent Sicherheit (1 bis 5):~\rule{2.5em}{0.4pt}\par
\end{minipage}

\medskip

\noindent\begin{minipage}{\linewidth}
\textbf{Frage 3}\par
Wird der Import der Studierenden noch an einer anderen Stelle der Anwendung benutzt?

\begin{itemize}[nosep,leftmargin=2.4em]
\item[\kreuz] a\quad Ja, auf der Seite mit den Team-Sets des Kurses
\item[\kreuz] b\quad Nein, nur auf der Seite mit den Studierenden
\item[\kreuz] c\quad Ja, auch im Einrichtungsassistenten des Kurses
\item[\kreuz] d\quad Weiß ich nicht
\end{itemize}
\noindent Sicherheit (1 bis 5):~\rule{2.5em}{0.4pt}\par
\end{minipage}

\medskip

\noindent\begin{minipage}{\linewidth}
\textbf{Frage 4}\par
Der Sections-Filter in der Leiste über der Tabelle: was passiert mit ihm, wenn die Spalte aus der Spaltenliste verschwindet?

\begin{itemize}[nosep,leftmargin=2.4em]
\item[\kreuz] a\quad Er bleibt stehen und zeigt dann eine leere Auswahl an
\item[\kreuz] b\quad Er verschwindet mit, weil er prüft, ob die Spalte da ist
\item[\kreuz] c\quad Er muss an seiner eigenen Stelle von Hand entfernt werden
\item[\kreuz] d\quad Weiß ich nicht
\end{itemize}
\noindent Sicherheit (1 bis 5):~\rule{2.5em}{0.4pt}\par
\end{minipage}

\medskip

\noindent\begin{minipage}{\linewidth}
\textbf{Frage 5}\par
Womit ist die Anfrage an den Server für die Studierendenliste gekennzeichnet?

\begin{itemize}[nosep,leftmargin=2.4em]
\item[\kreuz] a\quad Mit dem Kurs und den Einstellungen der Tabelle (Seite, Sortierung, Suche)
\item[\kreuz] b\quad Nur mit der Kennung des Kurses, die Einstellungen bleiben im Browser
\item[\kreuz] c\quad Mit der Kennung der angemeldeten Person und dem Kurs, sonst nichts
\item[\kreuz] d\quad Weiß ich nicht
\end{itemize}
\noindent Sicherheit (1 bis 5):~\rule{2.5em}{0.4pt}\par
\end{minipage}

\medskip

\noindent\begin{minipage}{\linewidth}
\textbf{Frage 6}\par
Wie viele Studierende zeigt die Tabelle beim ersten Öffnen je Seite?

\begin{itemize}[nosep,leftmargin=2.4em]
\item[\kreuz] a\quad 25
\item[\kreuz] b\quad 10
\item[\kreuz] c\quad 50
\item[\kreuz] d\quad Weiß ich nicht
\end{itemize}
\noindent Sicherheit (1 bis 5):~\rule{2.5em}{0.4pt}\par
\end{minipage}

\medskip

\noindent\begin{minipage}{\linewidth}
\textbf{Frage 7 (nur im Nachher-Bogen)}\par
Das View-Menü über der Tabelle: woher weiß es, welche Spalten es anbietet?

\begin{itemize}[nosep,leftmargin=2.4em]
\item[\kreuz] a\quad Aus einer eigenen Liste, die im Menü fest eingetragen ist
\item[\kreuz] b\quad Vom Server, der die Spalten zusammen mit der Liste schickt
\item[\kreuz] c\quad Aus der Spaltenliste der Tabelle, alle ausblendbaren Spalten
\item[\kreuz] d\quad Weiß ich nicht
\end{itemize}
\noindent Sicherheit (1 bis 5):~\rule{2.5em}{0.4pt}\par
\end{minipage}

\medskip

\noindent\begin{minipage}{\linewidth}
\textbf{Frage 8 (nur im Nachher-Bogen)}\par
Die Einstellungen, die die Tabelle gerade hat: wo werden sie gehalten?

\begin{itemize}[nosep,leftmargin=2.4em]
\item[\kreuz] a\quad In der Tabelle selbst, getrennt von den Daten
\item[\kreuz] b\quad Im selben Baustein, der auch die Liste holt
\item[\kreuz] c\quad Auf dem Server, der sie je Person speichert
\item[\kreuz] d\quad Weiß ich nicht
\end{itemize}
\noindent Sicherheit (1 bis 5):~\rule{2.5em}{0.4pt}\par
\end{minipage}

\medskip

\subsection*{Teil C: Was hat der Agent gemacht? (nur im Nachher-Bogen)}

\noindent\begin{minipage}{\linewidth}
\textbf{Frage 9}\par
An welchen dieser Stellen hat der Agent etwas geändert? (mehrere möglich)

\begin{itemize}[nosep,leftmargin=2.4em]
\item[\kreuz] a\quad In der Liste, die die Spalten festlegt
\item[\kreuz] b\quad Im Baustein, der die Liste holt und an die Tabelle gibt
\item[\kreuz] c\quad Auf der Seite, in der Kopfzeile mit dem Import-Knopf
\item[\kreuz] d\quad Im Baustein, der den Import ausführt
\item[\kreuz] e\quad Weiß ich nicht
\end{itemize}
\noindent Sicherheit (1 bis 5):~\rule{2.5em}{0.4pt}\par
\end{minipage}

\medskip

\noindent\begin{minipage}{\linewidth}
\textbf{Frage 10}\par
Hat der Agent sein Ergebnis geprüft, bevor er fertig war (Typprüfung, Lint)?

\begin{itemize}[nosep,leftmargin=2.4em]
\item[\kreuz] a\quad Ja
\item[\kreuz] b\quad Nein
\item[\kreuz] c\quad Weiß ich nicht
\end{itemize}
\end{minipage}

\medskip

\subsection*{Teil D: Wie war die Aufgabe für dich (nur im Nachher-Bogen)}

Je Zeile eine Zahl von 0 bis 100 in Fünferschritten.

\noindent Geistige Anforderung (wie viel Denken, Entscheiden, Suchen war nötig):~\rule{2.5em}{0.4pt}\par
\noindent Körperliche Anforderung:~\rule{2.5em}{0.4pt}\par
\noindent Zeitliche Anforderung (wie viel Zeitdruck hast du empfunden):~\rule{2.5em}{0.4pt}\par
\noindent Leistung (0 heißt sehr gut gelungen, 100 heißt misslungen):~\rule{2.5em}{0.4pt}\par
\noindent Anstrengung (wie hart musstest du arbeiten):~\rule{2.5em}{0.4pt}\par
\noindent Frustration (wie unsicher, entmutigt, gereizt warst du):~\rule{2.5em}{0.4pt}\par
Je Zeile eine Zahl von 1 (stimme nicht zu) bis 5 (stimme voll zu).

\noindent Ich wusste während des Laufs, woran der Agent gerade arbeitet:~\rule{2.5em}{0.4pt}\par
\noindent Ich könnte jemandem aus dem Team erklären, was geändert wurde und warum:~\rule{2.5em}{0.4pt}\par
\noindent In den Sitzungen trug jede Skala Ankertexte für die Werte 0 und 100. Eine hohe Zahl bedeutet auf jeder Skala eine höhere Beanspruchung.

\begin{center}
\small
\begin{tabular}{lll}
\toprule
Skala & 0 & 100 \\
\midrule
Geistige Anforderung & wenig & sehr viel \\
Körperliche Anforderung & keine & sehr hoch \\
Zeitliche Anforderung & keinen Zeitdruck & sehr starken Zeitdruck \\
Leistung & sehr gut gelungen & misslungen \\
Anstrengung & gar nicht & sehr hart \\
Frustration & gar nicht & sehr \\
\bottomrule
\end{tabular}
\end{center}
\subsection*{Teil A: Welche Stellen spielen eine Rolle?}

Welche dieser Stellen im Projekt spielen eine Rolle dabei, welche Spalten die Tabelle zeigt und wann sie ihre Daten neu lädt?

\noindent\begin{minipage}{\linewidth}
\textbf{Stelle 1}\par
Die Liste, die festlegt, welche Spalten die Tabelle hat

\begin{itemize}[nosep,leftmargin=2.4em]
\item[\kreuz] spielt eine Rolle
\item[\kreuz] spielt keine Rolle
\item[\kreuz] weiß ich nicht
\end{itemize}
\end{minipage}

\medskip

\noindent\begin{minipage}{\linewidth}
\textbf{Stelle 2}\par
Die Stelle, an der die Studierendenliste vom Server geholt wird

\begin{itemize}[nosep,leftmargin=2.4em]
\item[\kreuz] spielt eine Rolle
\item[\kreuz] spielt keine Rolle
\item[\kreuz] weiß ich nicht
\end{itemize}
\end{minipage}

\medskip

\noindent\begin{minipage}{\linewidth}
\textbf{Stelle 3}\par
Die Leiste über der Tabelle mit Suche, Filter und View-Menü

\begin{itemize}[nosep,leftmargin=2.4em]
\item[\kreuz] spielt eine Rolle
\item[\kreuz] spielt keine Rolle
\item[\kreuz] weiß ich nicht
\end{itemize}
\end{minipage}

\medskip

\noindent\begin{minipage}{\linewidth}
\textbf{Stelle 4}\par
Der Baustein, der den Import der Studierenden ausführt

\begin{itemize}[nosep,leftmargin=2.4em]
\item[\kreuz] spielt eine Rolle
\item[\kreuz] spielt keine Rolle
\item[\kreuz] weiß ich nicht
\end{itemize}
\end{minipage}

\medskip

\noindent\begin{minipage}{\linewidth}
\textbf{Stelle 5}\par
Die Meldungen unten rechts

\begin{itemize}[nosep,leftmargin=2.4em]
\item[\kreuz] spielt eine Rolle
\item[\kreuz] spielt keine Rolle
\item[\kreuz] weiß ich nicht
\end{itemize}
\end{minipage}

\medskip

\noindent\begin{minipage}{\linewidth}
\textbf{Stelle 6}\par
Die obere Navigation

\begin{itemize}[nosep,leftmargin=2.4em]
\item[\kreuz] spielt eine Rolle
\item[\kreuz] spielt keine Rolle
\item[\kreuz] weiß ich nicht
\end{itemize}
\end{minipage}

\medskip

\noindent\begin{minipage}{\linewidth}
\textbf{Stelle 7}\par
Die Seite mit den Team-Sets

\begin{itemize}[nosep,leftmargin=2.4em]
\item[\kreuz] spielt eine Rolle
\item[\kreuz] spielt keine Rolle
\item[\kreuz] weiß ich nicht
\end{itemize}
\end{minipage}

\medskip

\noindent\begin{minipage}{\linewidth}
\textbf{Stelle 8}\par
Die Anmeldung

\begin{itemize}[nosep,leftmargin=2.4em]
\item[\kreuz] spielt eine Rolle
\item[\kreuz] spielt keine Rolle
\item[\kreuz] weiß ich nicht
\end{itemize}
\end{minipage}

\medskip

